\documentclass[11pt]{article}
\usepackage[margin=1in]{geometry}
\usepackage[T1]{fontenc}
\usepackage[utf8]{inputenc}
\usepackage{lmodern,microtype,amsmath,amssymb,amsthm,mathtools}
\usepackage[colorlinks=true,linkcolor=blue,citecolor=blue,urlcolor=blue]{hyperref}
\hypersetup{pdftitle={Failure of Bayesian Gaussian process classification at exponentially large sample sizes},pdfauthor={ChatGPT-6 Astra; Kaarel Hanni; Dmitry Vaintrob}}
\setlength{\emergencystretch}{2em}
\newtheorem{theorem}{Theorem}[section]
\newtheorem{lemma}[theorem]{Lemma}
\newtheorem{proposition}[theorem]{Proposition}
\newtheorem{corollary}[theorem]{Corollary}
\theoremstyle{remark}\newtheorem{remark}[theorem]{Remark}
\DeclareMathOperator{\Var}{Var}
\DeclareMathOperator{\sign}{sign}
\DeclareMathOperator{\dist}{dist}
\title{Failure of Bayesian Gaussian process classification\\at exponentially large sample sizes}
\author{ChatGPT-6 Astra \and Kaarel H\"anni \and Dmitry Vaintrob\thanks{Authors are listed in alphabetical order.}}
\date{October 15, 2026}
\begin{document}\maketitle
\begin{abstract}
We study the exact Gaussian sign posterior for the saturated neural-network
Gaussian process on the Boolean cube. At every fixed depth, we prove that
there is an exponentially large sample schedule on which learning
full parity still leaves a fixed positive fraction of inputs without
correct-label probability greater than $2/3$, with probability at least
$3/4$ over independent uniform training inputs. The result also applies
to the width limit of globally margin-conditioned Bayesian $\tanh$
networks at every fixed scaling exponent $\gamma<0$. The proof combines
sharp Fourier level masses, a random tail-covariance operator bound,
and uniform ReLU sampling with exact constraints on the posterior mean.
It concerns the posterior predictive probability, rather than the sign
of a regression mean. The conclusion is failure at specified exponential
sample sizes; a lower bound covering every smaller sample size is not
proved. The fixed-margin case $\gamma=0$ and the regime $0<\gamma<1$
are outside the exponential result established here.
\end{abstract}

\section{Model, result, and scope}
Let $\mathcal X=\{-1,1\}^d$ with the uniform probability measure, and
let $f(x)=\prod_{j=1}^d x_j$ be full parity. Although parity has Boolean
circuits of size $O(d)$, its learnability depends on the specified prior
and prediction rule. We analyze Gaussian-process classification with
covariance
\begin{equation}\label{eq:paper-kernel}
 K(x,y)=H^{\circ(L-1)}\!\left(\frac{1+x\cdot y}{d+1}\right),
 \qquad H(t)=\frac2\pi\arcsin t,
\end{equation}
where the depth $L\ge2$ is fixed. For a centered Gaussian vector
$G\sim N(0,K)$ and training inputs $X_1,\ldots,X_n$, the posterior
conditions on the exact signs $f(X_i)G(X_i)>0$. Its probability of
being correct at $x$ is
\[
 q_S(x)=\mathbb P\{f(x)G(x)>0\mid f(X_i)G(X_i)>0\text{ for all }i\}.
\]
We count an input as a failure when $q_S(x)\le2/3$. Thus the criterion
concerns the actual posterior probability and includes equality as a
failure. It is not identified with the error of the sign of a posterior
mean.

Our main theorem, Theorem~\ref{thm:exponential-schedule}, gives constants
$A,\alpha,\varepsilon>0$ depending at most on the fixed depth. With
$r_d$ odd and $r_d/d\to\alpha$, put
\begin{equation}\label{eq:paper-schedule}
 D_{r_d}=\sum_{j=0}^{r_d}\binom dj,
 \qquad n_d=\left\lceil A D_{r_d}\log(D_{r_d}+1)\right\rceil.
\end{equation}
Then $n_d=\exp((h(\alpha)+o(1))d)$, where $h$ is the binary entropy
with natural logarithms. With probability at least $3/4$ over $n_d$
independent uniform training inputs, a fraction at least $\varepsilon$
of test inputs has correct-label probability at most $13/20<2/3$.
One conservative choice in the proof is
\[
 \varepsilon=\frac{2\pi(1/e-7/20)^2}{2048}>9\cdot10^{-7}.
\]
The constants are not optimized. In particular the theorem establishes
a fixed positive failure fraction, not a numerical failure fraction
such as $1/10$.

The quantifiers in this conclusion matter. It establishes failure on
the displayed exponential sample schedule and does not assert failure
at every smaller sample size. Posterior confidence for a fixed target
has not been shown to be monotone in the sample size. The argument also
does not establish an exponential statement at zero or positive
subcritical scaling exponents. The later paper~\cite{gaussian-all-n} proves failure at every specified
sample size below an exponential threshold, and
\cite{subcritical-exponential} also treats zero and positive subcritical
exponents. This paper retains the earlier schedule argument and its
quantifiers.

\subsection{The finite-width network and its Gaussian limit}
The same theorem applies to the following network convention. All
hidden layers have width $N$, with
\[
 h^0(x)=(1,x),\quad
 h^s(x)=\tanh(W^s h^{s-1}(x))\quad(1\le s<L),\quad
 z_N(x)=W^L h^{L-1}(x).
\]
There are no additional biases. The weights are independent centered
Gaussians with variances
\begin{equation}\label{eq:paper-weights}
 \Var(W^1_{ij})=N^{1-\gamma}/d,\quad
 \Var(W^s_{ij})=N^{-\gamma}\ (1<s<L),\quad
 \Var(W^L_i)=N^{-\gamma-1}.
\end{equation}
Equivalently, prescribe the normalized weight scale
$t(N)=N^{(1-\gamma)/2}$ before varying width. The three variances are
$t(N)^2/d$, $t(N)^2/N$, and $t(N)^2/N^2$, respectively.
Thus $\gamma$ is fixed while the scalar scale follows this width schedule.

Before training, condition on the global margin event
$A_a=\{|z_N(x)|>a\text{ for every }x\in\mathcal X\}$, where $a>0$
is fixed. Training then conditions on $\sign z_N(X_i)=f(X_i)$.
Let $q_{N,S}(x)$ be the resulting correct-label probability and
$F_N(S)=\mathbb P_x\{q_{N,S}(x)\le2/3\}$.

\begin{proposition}[Conditional width limit]\label{prop:paper-width}
At fixed $d,L,a,S$ and $\gamma<0$, for every $x\in\mathcal X$,
\[
 q_{N,S}(x)\longrightarrow q_S(x).
\]
All these posteriors are well defined for sufficiently large $N$.
Consequently
\[
 \liminf_{N\to\infty}F_N(S)\ge
 \mathbb P_x\{q_S(x)\le13/20\}.
\]
\end{proposition}
\begin{proof}
Let $Q_N^s(x,y)=N^{-1}\sum_i h_i^s(x)h_i^s(y)$. Conditional on
the preceding layer, hidden rows are independent bounded random vectors.
The conditional law of large numbers therefore applies to all empirical
second moments, since the cube is finite. Preactivation variances grow
like $N^{1-\gamma}$. Thus $\tanh$ converges to the sign function under
a Gaussian coupling. The identity
$\mathbb E[\sign U\sign V]=(2/\pi)\arcsin\rho$ for standard jointly
Gaussian $U,V$ of correlation $\rho$ follows, for example, by measuring
the corresponding angular sectors of a rotationally invariant planar
Gaussian. Dominated convergence and induction over layers yield
\[
 Q_N^{L-1}\longrightarrow K\quad\text{in probability}.
\]
For completeness, the convergence of the conditional moment map is
uniform near its limiting covariance: if it were not, choose a
converging covariance subsequence, couple by continuous matrix square
roots, and apply dominated convergence. Its diagonal is positive, so
the limiting scalar Gaussian coordinates vanish with probability zero.
The first layer normalizes $(d+1)/d$ on the diagonal and hence gives
the argument $(1+x\cdot y)/(d+1)$ in \eqref{eq:paper-kernel}.

Conditionally on the hidden layers,
$N^{\gamma/2}z_N\sim N(0,Q_N^{L-1})$. Therefore the whole output
vector $G_N=N^{\gamma/2}z_N$ converges in law to $G\sim N(0,K)$.
The global margin in these coordinates is
$|G_N(x)|>aN^{\gamma/2}$, whose threshold tends to zero. The spectral
proof below shows $K$ is positive definite. Hence all coordinate
hyperplanes are Gaussian null sets and every sign orthant has positive
probability. A shrinking-neighborhood argument around the coordinate
hyperplanes shows that imposing this vanishing margin changes the
probability of each fixed sign event by $o(1)$. Ratios of probabilities
of sign events now give $q_{N,S}(x)\to q_S(x)$ and positivity of the
conditioning probability for large $N$.

At a point where $q_S(x)\le13/20$, convergence implies
$q_{N,S}(x)<2/3$ eventually. The cube is finite, so the implication
holds simultaneously at all these points, proving the final assertion.
\end{proof}

The width limit is always taken at fixed $d,L,\gamma,a,S$. The
subsequent exponential statement concerns $d\to\infty$ for these
limiting learners. No uniform interchange of the width and dimension
limits is asserted.

\paragraph{Proof roadmap.}
The Fourier estimates control both the trace and largest eigenvalue of
the high-degree covariance. Random-submatrix bounds then control that
covariance on the sample. Convex Gaussian sign conditioning bounds the
actual posterior mean and turns confident predictions into latent
margins. Uniform sampling of its low-degree part completes the
contradiction at the specified sample schedule.

\section{Fourier level masses of the kernel}
Put $D=d+1$, $\beta=2^{-(L-1)}$, and $F=H^{\circ(L-1)}$.
The characters $\chi_U(x)=\prod_{j\in U}x_j$ diagonalize the
normalized covariance operator $M^{-1}K$, $M=2^d$. Let $\lambda_j$
be its eigenvalue at degree $j$. We first work on the full $D$-cube
with kernel $F(u\cdot v/D)$, whose degree-$j$ eigenvalue is denoted
by $a_j$. Constants in this and the next section may depend on $L$.
Notation $b\asymp c$ means two inequalities by positive constant
multiples, uniformly in the stated ranges.

\begin{lemma}[Power coefficients and paired eigenvalues]\label{lem:power}
Write $F(z)=\sum_{m\text{ odd}}c_mz^m$. Every odd coefficient is
positive, $\sum_m c_m=1$, and
\[
 c_m\asymp\frac{\Gamma(m-\beta)}{m!}\asymp m^{-1-\beta}
 \qquad(m\ge1\text{ odd}).
\]
The full-cube eigenvalues vanish at even degrees and strictly decrease
through odd degrees. They are positive at every odd degree. On the
original $d$-cube,
\[
 \lambda_{2k}=\lambda_{2k+1}=a_{2k+1}>0,
\]
with a final unpaired eigenvalue if $d$ is even. In particular $K$
is positive definite.
\end{lemma}
\begin{proof}
The series of $H$ has positive coefficients at every odd exponent,
with
\[
 [z^{2k+1}]H(z)=\frac2\pi\frac{\binom{2k}{k}}{4^k(2k+1)}.
\]
Composition preserves strict positivity and oddness. Near $z=1$ in
a slit neighborhood,
\[
 H(z)=1-b(1-z)^{1/2}+O((1-z)^{3/2}),
 \qquad b=\frac{2\sqrt2}{\pi}.
\]
Induction gives
$F(z)=1-A(1-z)^\beta+O((1-z)^{3\beta})$, with $A>0$;
at each composition its leading constant changes from $A_t$ to
$b\sqrt{A_t}$, starting from $A_1=b$.
These expansions admit algebraic coefficient transfer
\cite[Chapter VI]{flajolet2009}. Here are the analytic-domain details.
On the unit circle outside $\{-1,1\}$, positivity of every odd
coefficient gives $|H(z)|<1$: equality in the triangle inequality
would require $z^2=1$. Hence the iterated functions extend
analytically through every compact arc away from these endpoints.
At $1$, restrict $\arg(1-z)$ to a closed proper subinterval of
$(-\pi,\pi)$. Each square-root composition reduces this angular
range, and its positive leading coefficient keeps its image in the
slit neighborhood for the next arcsine composition. Oddness gives
the identical argument at $-1$. These local sectors and the compact
arcs provide the required two-singularity indented domain.
The contributions of $1$ and $-1$ add at odd exponents, giving
\[
 c_m\sim\frac{2A}{-\Gamma(-\beta)}m^{-1-\beta}
 \quad(m\to\infty\text{ odd}).
\]
Together with positivity of the finitely many initial coefficients
and the gamma-ratio asymptotic, this proves the coefficient comparisons.
For $L=2$ the explicit formula and Stirling's bounds suffice instead.

Let $p_m(j)$ be the probability that the simple coordinate-flip walk
on the $D$-cube reaches a specified distance-$j$ vertex in $m$ steps.
Expanding $(u\cdot v/D)^m$ shows that its degree-$j$ Fourier coefficient
is $p_m(j)$. Thus $a_j=\sum_{m\text{ odd}}c_mp_m(j)$.
For the walk recursion,
\begin{align*}
 D[p_{m+1}(j)-p_{m+1}(j+2)]
 ={}&j[p_m(j-1)-p_m(j+1)]\\
 &+(D-j-2)[p_m(j+1)-p_m(j+3)].
\end{align*}
Induction from time zero gives nonincreasing probabilities through
each parity class; terms with zero boundary coefficients are omitted.
The term $m=j$ makes the odd-eigenvalue inequalities strict and proves
positivity. Even-degree eigenvalues vanish by parity.
Restricting the distinguished coordinate to $+1$ gives
$\lambda_q=a_q+a_{q+1}$, proving the claims.
\end{proof}

For odd $r$, the exact combined multiplicity of reduced degrees
$j-1,j$ is $\binom d{j-1}+\binom dj=\binom Dj$. Consequently
\begin{equation}\label{eq:paired}
 D_r=\sum_{\substack{j\le r\\j\text{ odd}}}\binom Dj,
 \quad
 \tau_r=\sum_{\substack{j>r\\j\text{ odd}}}\binom Dj a_j,
 \quad \lambda_{r+1}=a_{r+2}.
\end{equation}
The odd cutoff is essential for these exact identities.

\begin{lemma}[Sharp level and tail estimates]\label{lem:sharp}
Uniformly for odd $1\le j\le D/4$ and odd $1\le r\le D/8$,
\begin{align}
 \binom Dj a_j&\asymp j^{-1-\beta},\label{eq:level}\\
 \tau_r&\asymp r^{-\beta},\label{eq:trace}\\
 \frac{\lambda_{r+1}}{\tau_r}
 &\asymp\frac1{r\binom D{r+2}},\label{eq:ratio}\\
 s_{2,r}:=\sum_{q>r}\binom dq\lambda_q^2
 &\le\frac{C r^{-2\beta-2}}{\binom D{r+2}}.
 \label{eq:secondmoment}
\end{align}
The same level-mass estimate holds up to $j\le\rho D$ for any fixed
$\rho<1/2$, with constants depending on $\rho$.
\end{lemma}
\begin{proof}
The gamma integral and Lemma~\ref{lem:power} give, for odd $j$,
\begin{align}
 a_j&\asymp D^{-\beta}\int_0^\infty
 u^{-1-\beta}p(u)^j(1-p(u))^{D-j}\,du,
 \quad p(u)=\frac{1-e^{-2u}}2,\label{eq:poisson}\\
 &\asymp D^{-\beta}\int_0^{1/2}p^j(1-p)^{D-j}w_\beta(p)\,dp,
 \label{eq:betaweight}
\end{align}
where $w_\beta(p)=[-\frac12\log(1-2p)]^{-1-\beta}/(1-2p)$.
Indeed the gamma integral Poissonizes the walk time, and independent
Poisson coordinate counts give independent coordinate parities with
success probability $p(u)$. Since an odd-distance vertex has zero
probability at even walk times, no extra parity factor occurs.
The weight $w_\beta$ is integrable near $1/2$ and comparable to
$p^{-1-\beta}$ on $[0,b]$ for every fixed $b<1/2$.

Fix $\rho<b<1/2$ and take $j\le\rho D$. The integral on $[0,b]$
is comparable to the complete beta integral
$B(j-\beta,D-j+1)$. For the lower comparison, a beta random variable
with these parameters has mean at most $\rho$, so Markov's inequality
puts at least a fraction $1-\rho/b$ of its mass in $[0,b]$.
The identity
\[
 \binom Dj B(j-\beta,D-j+1)
 =\frac{\Gamma(D+1)\Gamma(j-\beta)}
        {\Gamma(j+1)\Gamma(D+1-\beta)}
 \asymp D^\beta j^{-1-\beta}
\]
gives the desired mass. On $[b,1/2]$, the binomial factor decreases
with $p$ and, after multiplication by $\binom Dj$, is at most
$\exp[-D\operatorname{KL}(\rho\Vert b)]$. The remaining integrable
weight makes this contribution exponentially small. This proves
\eqref{eq:level} and its stated extended range.

Sum \eqref{eq:poisson} over the odd distances above $r$. Dropping the
parity restriction and using the binomial Markov bound gives
\[
 \tau_r\le C D^{-\beta}\int_0^\infty
 u^{-1-\beta}\min\{1,Du/(r+1)\}\,du
 \le C(r+1)^{-\beta}.
\]
For the lower bound, sum \eqref{eq:level} over the odd $j$ in
$(r,2r+2]$, using a fixed cutoff strictly below $D/2$ as needed at
the endpoint. This yields $c r^{-\beta}$; bounded $r$ are absorbed
by a change of constants. Equation~\eqref{eq:ratio} follows from
\eqref{eq:level}, \eqref{eq:trace}, and \eqref{eq:paired}.

For the second moment, use paired multiplicities to write
$s_{2,r}=\sum_{j>r,\ j\text{ odd}}\binom Dj a_j^2$.
Up to an odd cutoff $b_D$ near $D/4$, its terms are at most
$Cj^{-2\beta-2}/\binom Dj$. Consecutive odd inverse binomial
coefficients have ratio
\[
 \frac{(j+1)(j+2)}{(D-j)(D-j-1)},
\]
uniformly less than a fixed constant below one in this range.
The sum is therefore bounded by a constant times its first term,
at $j=r+2$. Beyond $b_D$, monotonicity gives the bound
$a_{b_D}\tau_{b_D}\le C D^{-2\beta-1}/\binom D{b_D}$,
which is exponentially smaller than the claimed bound uniformly
for $r\le D/8$. This proves \eqref{eq:secondmoment}.
\end{proof}

Finally, two-step binomial ratios give, uniformly for odd $r\le D/8$,
\begin{equation}\label{eq:dimensionratios}
 \binom Dr\le D_r\le C\binom Dr,
 \qquad \frac{D_r}{\binom D{r+2}}\le C(r/D)^2.
\end{equation}
All constants used here can be fixed independently of a subsequently
chosen small positive limit $r/d\to\alpha$.

\section{Random tail covariance matrices}
Let $Q$ be the covariance kernel retaining only degrees greater than an
odd cutoff $r$, and write $\tau=\tau_r$, $\ell=\lambda_{r+1}$, and
$s_2=s_{2,r}$. This section proves the sample event needed later, without
assuming that a Gaussian approximation remains valid after conditioning.
The entire argument concerns the exact limiting Gaussian covariance.

\begin{lemma}[Far-pair decay]\label{lem:far}
There are constants $c_0,C>0$ such that, when $r\le D/8$ is odd,
\[
 \max_{d/4\le\dist(x,y)\le3d/4}|Q(x,y)|
 \le\rho_r:=C\sqrt D\,e^{-r/8}+Ce^{-c_0D}.
\]
\end{lemma}
\begin{proof}
The full-cube Krawtchouk polynomial satisfies
\[
 \sum_{j=0}^D\mathsf K_j(h;D)z^j=(1+z)^{D-h}(1-z)^h.
\]
If $\delta D\le h\le(1-\delta)D$, put $b=\min(h,D-h)$.
On $|z|=t$ the generating function is bounded in absolute value by
\[
 (1+t^2)^b(1+t)^{D-2b}
 \le(1+t)^D\exp\!\left[-\frac{2\delta Dt}{(1+t)^2}\right].
\]
Cauchy's coefficient bound at $t=j/(D-j)$ and the uniform Stirling
lower bound for $\binom Dj$ therefore give, for $1\le j<D$,
\begin{equation}\label{eq:kraw}
 \frac{|\mathsf K_j(h;D)|}{\binom Dj}
 \le C\sqrt D\exp[-2\delta j(1-j/D)].
\end{equation}
For reduced-cube points at distance in $[d/4,3d/4]$, their augmented
full-cube distance lies in $[D/8,7D/8]$ for large $d$.
Because $r$ is odd, the exact tail is
$Q(x,y)=\sum_{j>r,\ j\text{ odd}}a_j\mathsf K_j(h;D)$.
For $r<j\le D-r$, \eqref{eq:kraw} gives a bound
$C\sqrt D e^{-r/8}$ after summation, since the total spectral mass
is one. For $j>D-r$, monotonicity and an odd cutoff $b_D$ near $D/4$
give
\[
 \sum_{j>D-r}\binom Dj a_j
 \le a_{b_D}\sum_{k<r}\binom Dk
 \le C D^{-\beta-1/2}
       \exp[-D(h(1/4)-h(1/8))+O(1)].
\]
The positive entropy gap absorbs the polynomial factor and proves the
claim. Adjust constants for bounded $d$.
\end{proof}

\begin{lemma}[A tail Gram deviation bound]\label{lem:gram-general}
Let $X_1,\ldots,X_n$ be independent uniform cube points, with
$1\ll n\le M/4$. Let
\[
 B_{\mathrm{near}}=\#\{y:\dist(x,y)\notin[d/4,3d/4]\}
 \le 2e^{d h(1/4)},\qquad q=2\log M.
\]
The count is independent of $x$. For every $z>0$,
\begin{align}
 &\mathbb P\{\|Q_{SS}-\tau I_n\|_{\mathrm{op}}>z\tau\}
 \notag\\
 &\quad\le\frac C{z^2}
 \left[\frac{q\rho_r}{\tau}
       +\frac{\sqrt{nq s_2}}{\tau}
       +\frac{n\ell}{\tau}
       +\frac{n(1+B_{\mathrm{near}})}M\right]^2
       +\frac{n^2(1+B_{\mathrm{near}})}M.
 \label{eq:gram-bound}
\end{align}
Here the sampled matrix includes repeated training locations.
\end{lemma}
\begin{proof}
Let $J=Q-\tau I_M$, a zero-diagonal Hermitian matrix. Form a deterministic
mask $J_{\mathrm{far}}$ retaining only pairs at distances in
$[d/4,3d/4]$. Since $Q$ is positive semidefinite with constant diagonal
$\tau$, every entry has absolute value at most $\tau$. Thus
\begin{align*}
 \|J_{\mathrm{far}}\|_{\max}&\le\rho_r,\\
 \|J_{\mathrm{far}}\|_{1\to2}^2&\le M s_2,\\
 \|J_{\mathrm{far}}\|_{\mathrm{op}}
 &\le M\ell+\tau+\tau B_{\mathrm{near}}.
\end{align*}
For the second inequality, Parseval gives squared column norm
$M s_2-\tau^2$ for $J$, and entrywise masking only decreases it.
For the third, $\|J\|\le M\ell+\tau$, while the removed matrix has
maximum absolute row sum at most $\tau B_{\mathrm{near}}$.
The entropy bound on $B_{\mathrm{near}}$ is the elementary binomial
tail bound on distances to a fixed cube point.

Let $R$ retain each of the $M$ coordinates independently with
probability $p=2n/M$. The zero-diagonal case of Tropp's random
principal-submatrix theorem \cite[Theorem 1.1]{tropp2008}, at moment
order two and $q=2\log M$, states
\[
 (\mathbb E\|RJR\|^2)^{1/2}
 \le C\left[q(\mathbb E\|RJR\|_{\max}^2)^{1/2}
       +\sqrt{pq}(\mathbb E\|JR\|_{1\to2}^2)^{1/2}
       +p\|J\|\right].
\]
Apply it to $J_{\mathrm{far}}$ and substitute the preceding bounds:
\begin{equation}\label{eq:bernoulli-gram}
 (\mathbb E\|R J_{\mathrm{far}}R\|^2)^{1/2}
 \le C[q\rho_r+\sqrt{nq s_2}+n\ell
               +(n/M)\tau(1+B_{\mathrm{near}})].
\end{equation}
The logarithm here is that of the original dimension $M$, not the
sample size.

To pass to a uniform subset of size $n$, use a uniform permutation of
all cube points, independently of $B\sim\operatorname{Bin}(M,2n/M)$.
Its first $B$ elements form the Bernoulli subset. On $B\ge n$, the
first $n$ principal submatrix is a compression of the larger one, so
its operator norm cannot be larger. The size-$n$ event is independent
of $B$, and $\mathbb P(B\ge n)\ge1/2$ for large $n$. Thus Markov's
inequality applied to \eqref{eq:bernoulli-gram} gives the same fixed-size
probability bound with a constant factor loss.

Conditional on no collisions, an independent size-$n$ sample is a
uniform size-$n$ subset, up to an irrelevant ordering. Collisions have
probability at most $n(n-1)/(2M)$. The probability of any sampled pair
violating the distance condition is at most $n^2B_{\mathrm{near}}/M$.
On their complement, the sampled matrices $J$ and $J_{\mathrm{far}}$
agree. Adding these exceptional probabilities proves
\eqref{eq:gram-bound}.
\end{proof}

\begin{corollary}[Simultaneous spectral and Gram control]
\label{cor:gram-schedule}
Fix $A>0$ and arbitrary constants $\eta,\xi>0$. There is a sufficiently
small $\alpha>0$ such that, whenever $r_d$ is odd with $r_d/d\to\alpha$
and $n_d=\lceil A D_{r_d}\log(D_{r_d}+1)\rceil$, for all sufficiently
large $d$,
\[
 \frac{n_d\lambda_{r_d+1}}{\tau_{r_d}}\le\xi,
 \qquad
 \mathbb P_S\{\|Q_{SS}\|_{\mathrm{op}}>2\tau_{r_d}\}\le\eta.
\]
\end{corollary}
\begin{proof}
Equations~\eqref{eq:ratio}, \eqref{eq:secondmoment}, and
\eqref{eq:dimensionratios} give
\begin{align*}
 \frac{n_d\ell}{\tau}
 &\le C_L A\alpha^2\log(e/\alpha)+o(1),\\
 \frac{n_dq s_2}{\tau^2}
 &\le C_L A\alpha\log(e/\alpha)+o(1).
\end{align*}
Also $q\rho_r/\tau$ tends to zero exponentially. Choose $\alpha$
small enough that $2h(\alpha)<\log2-h(1/4)$. Then all near-pair and
collision terms in \eqref{eq:gram-bound} tend to zero exponentially,
because $n_d=\exp((h(\alpha)+o(1))d)$. The remaining bracket in
\eqref{eq:gram-bound} is at most
\[
 C_L\left[\sqrt{A\alpha\log(e/\alpha)}
                 +A\alpha^2\log(e/\alpha)\right]+o(1).
\]
It tends to zero as $\alpha\downarrow0$ after $A$ is fixed. Taking
$z=1$ in that bound and then choosing $\alpha$ still smaller proves
both assertions.
\end{proof}

% Gaussian-posterior part of the exponential sample-schedule theorem.
% Requires theorem, lemma, proposition, corollary environments and amsmath,
% amssymb, amsthm. The spectral/random-Gram lemma is supplied separately.

\section{The actual Gaussian sign posterior}
Let $\mathcal X=\{-1,1\}^d$, $M=2^d$, and
$f(x)=\prod_{j=1}^d x_j$. All function norms and expectations on
$\mathcal X$ use its uniform probability measure. Let $G$ be a
nondegenerate centered Gaussian vector indexed by $\mathcal X$, with
covariance matrix $K$. The training inputs $X_1,\ldots,X_n$ are independent
and uniform. We condition on
\[
 \mathcal A_S=\{f(X_i)G(X_i)>0:1\le i\le n\}.
\]
The prediction at $x$ is its actual conditional correct-label probability
$q_S(x)=\mathbb P(f(x)G(x)>0\mid\mathcal A_S)$, and we write
\[
 F_\theta(S)=\mathbb P_x\{q_S(x)\le\theta\},\qquad
 \theta=\frac{13}{20},\qquad
 c=\sqrt{2\pi}\left(\frac1e-\frac7{20}\right)>0.
\]
The threshold $\theta<2/3$ supplies a strict buffer for passage from
limiting Gaussian posteriors to finite-width posterior probabilities.

Suppose the Walsh characters are eigenfunctions of the normalized
covariance operator, with degree-$j$ eigenvalue $\lambda_j$. Write
$V_r=\operatorname{span}\{\chi_U:|U|\le r\}$ and $D_r=\dim V_r$.
For $r<d$, let $Q$ be the covariance contributed by the complementary
Fourier modes, and put
\[
 \tau=Q(x,x)=\sum_{j>r}\binom dj\lambda_j,
 \qquad \ell=\max_{j>r}\lambda_j.
\]
The matrix eigenvalues of $Q$ are $M\lambda_j$; this normalization is
used throughout. Define $\|h\|_K^2=h^TK^{-1}h$, and write $u_+=\max(u,0)$ and $u_-=\max(-u,0)$.

\begin{lemma}[Mean norm and a confidence margin]\label{lem:schedule-mean}
Let $m(x)=\mathbb E[G(x)\mid\mathcal A_S]$, and let $s\le n$ be the
number of distinct sampled inputs. Then
\[
 \|m\|_K^2\le s,\qquad f(X_i)m(X_i)>0.
\]
Let $\sigma_x^2=\operatorname{Var}(G(x)\mid G_S)$, where $S$ here lists
the distinct inputs. Whenever $\sigma_x>0$ and $q_S(x)>\theta$,
\[
 f(x)m(x)>c\sigma_x.
\]
\end{lemma}
\begin{proof}
Whiten the Gaussian training vector as $G_S=K_{SS}^{1/2}\xi$.
The sign constraints form a homogeneous cone in $\xi$. Conditioning
leaves the squared Gaussian radius with mean $s$, so Jensen's inequality
and Gaussian regression give
\[
 m_S^TK_{SS}^{-1}m_S\le s,\qquad
 m=K_{\mathcal X,S}K_{SS}^{-1}m_S.
\]
The RKHS norm of this regression vector is the preceding quadratic
form. Its observed signed coordinates are strictly positive.

For the confidence claim, the scalar conditional variable
$Y=f(x)G(x)$ has a log-concave density: it is a marginal of a Gaussian
restricted to a convex set, by the marginal log-concavity theorem \cite{prekopa1973}. If $\mu=\mathbb E Y$ and $F_Y$ is its
cdf, log-concavity and Jensen's inequality imply
\[
 \log F_Y(\mu)\ge\mathbb E\log F_Y(Y)=-1.
\]
Consequently $\mathbb P(Y>\mu)\le1-e^{-1}<\theta$.
Conditioning first on the full training vector shows that the density
of $Y$ is a mixture of Gaussian densities with common variance
$\sigma_x^2$, and is therefore at most
$(\sqrt{2\pi}\sigma_x)^{-1}$. If $q_S(x)>\theta$, then $\mu>0$ and
\[
 \theta<\mathbb P(Y>0)
 \le 1-e^{-1}+\frac{\mu}{\sqrt{2\pi}\sigma_x}.
\]
This gives the stated margin.
\end{proof}

\begin{lemma}[Most residual variances remain large]\label{lem:schedule-var}
For every training sample,
\[
 \mathbb P_x\{\sigma_x^2<\tau/2\}\le \frac{2n\ell}{\tau}.
\]
Also, if $p$ is the degree-$r$ Fourier truncation of $m$ and $t=m-p$,
then $\|t\|_2^2\le n\ell$.
\end{lemma}
\begin{proof}
Write $G=G_{\le r}+G_{>r}$ as a sum of independent fields. Conditioning
additionally on the full low-degree field can only decrease conditional
Gaussian variance. The remaining covariance is
\[
 Q-Q_{\mathcal X,S}Q_{SS}^{\dagger}Q_{S,\mathcal X}.
\]
Its subtracted term is positive semidefinite, bounded above by $Q$,
and has rank at most $s\le n$. Its normalized trace is at most $n\ell$.
Markov's inequality proves the first assertion; it includes observed
points automatically. Fourier orthogonality and the mean norm bound
prove $\|t\|_2^2\le\ell\|m\|_K^2\le n\ell$.
\end{proof}

\section{Uniform sampling detects the negative part}
\begin{lemma}[Uniform ReLU sampling on the cube]\label{lem:schedule-relu}
For every fixed $\delta,\eta\in(0,1)$, there is a universal constant
$C$ such that, if
\[
 n\ge C\eta^{-2}D_r\log\!\left(\frac{D_r}{\delta\eta}\right),
\]
then with probability at least $1-\delta$, simultaneously for all
$p\in V_r$,
\[
 \left|\frac1n\sum_{i=1}^n(-f(X_i)p(X_i))_+
             -\mathbb E_x(-f(x)p(x))_+\right|
 \le\eta\mathbb E_x|p(x)|.
\]
In particular, when $r<d$ and $\eta=1/4$, the empirical negative
part is at least $\mathbb E|p|/4$.
\end{lemma}
\begin{proof}
Let $A$ have rows $f(x)(\chi_U(x))_{|U|\le r}$, one for each cube
point. Then $A^TA=MI$ and every row has squared norm $D_r$.
The $\ell_1$ Lewis weights of $A$ are exactly $D_r/M$: substitution
in their defining equation
\[
 w_x^2=A_x(A^T\operatorname{diag}(w^{-1})A)^{-1}A_x^T
\]
verifies this constant solution, and Lewis weights are unique.
The additive ReLU coreset theorem of Mai, Musco, and Rao
\cite[Theorem~5]{mai2021coresets} therefore samples precisely the uniform
distribution on the cube and gives the stated estimate. Its sample bound
is $O(\eta^{-2}D_r\log(D_r/(\delta\eta)))$; no additional logarithmic
factor is required. Finally $\mathbb E fp=0$, because $r<d$, so
$\mathbb E(-fp)_+=\mathbb E|p|/2$.
\end{proof}

\begin{proposition}[A deterministic reduction to a tail Gram bound]
\label{prop:schedule-deterministic}
Suppose a design obeys the conclusion of Lemma~\ref{lem:schedule-relu}
with $\eta=1/4$, and also
\begin{align}
 \|Q_{SS}\|_{\mathrm{op}}&\le B\tau,\label{eq:schedule-gram}\\
 \frac1n\#\{i:f(X_i)p(X_i)<0\}
       &\le\mathbb P_x\{f(x)p(x)\le0\}+u
       \quad\hbox{for all }p\in V_r.\label{eq:schedule-vc}
\end{align}
Here $B\ge1$, and the matrix in \eqref{eq:schedule-gram} lists the training sample
with multiplicity. Put $\zeta=n\ell/\tau$. If
\[
 u\le\frac{c^2}{2048B},\qquad
 \left(2+\frac8{c^2}\right)\zeta\le\frac{c^2}{2048B},
\]
then the actual Gaussian sign posterior satisfies
\[
 F_\theta(S)\ge\frac{c^2}{1024B}.
\]
\end{proposition}
\begin{proof}
Let $m=p+t$ be the low/high decomposition of its conditional mean,
and write $E_p=\mathbb P_x\{fp\le0\}$. Training consistency implies
\[
 (f(X_i)p(X_i))_-\le|t(X_i)|\mathbf1_{\{f(X_i)p(X_i)<0\}}.
\]
On the tail subspace $\|t\|_Q=\|t\|_K$. RKHS evaluation and $\|t\|_K^2\le n$ give
\[
 \frac1n\sum_i t(X_i)^2
 \le\frac{\|t\|_K^2}{n}\|Q_{SS}\|_{\mathrm{op}}\le B\tau.
\]
Combining the ReLU estimate, Cauchy--Schwarz, and
\eqref{eq:schedule-vc} yields
\begin{equation}\label{eq:schedule-l1upper}
 \mathbb E|p|\le4\sqrt{B\tau(E_p+u)}.
\end{equation}

Lemma~\ref{lem:schedule-var} and Markov's inequality show that, outside
sets of total measure at most
$a_0=(2+8/c^2)\zeta$, both
$\sigma_x\ge\sqrt{\tau/2}$ and
$|t(x)|\le c\sqrt\tau/(2\sqrt2)$ hold. At every remaining input
with confidence greater than $\theta$, Lemma~\ref{lem:schedule-mean}
gives
\[
 f(x)p(x)>\frac{c\sqrt\tau}{2\sqrt2}.
\]
Writing $e=F_\theta(S)+a_0$, we obtain $E_p\le e$ and
\[
 \frac{\mathbb E|p|}{\sqrt\tau}
 \ge\frac c{2\sqrt2}(1-e).
\]
If $e\le1/2$, comparison with \eqref{eq:schedule-l1upper} gives
$e+u\ge c^2/(512B)$, and hence
$F_\theta(S)\ge c^2/(1024B)$ under the stated assumptions.
If $e>1/2$, the same claimed lower bound follows directly from the
upper bound on $a_0$.
\end{proof}

\section{An exponentially large sample schedule}
The preceding proposition is deterministic. The remaining ingredients
are a standard VC estimate and the spectral/random-Gram estimates proved
in the preceding sections. We state their exact required form to make
the scope of the conclusion explicit.

\begin{theorem}[Failure along an exponential sample schedule]
\label{thm:exponential-schedule}
Fix the depth $L\ge2$. For the covariance
\[
 K(x,y)=H^{\circ(L-1)}\!\left(\frac{1+x\cdot y}{d+1}\right),
 \qquad H(t)=\frac2\pi\arcsin t,
\]
there are constants $A>0$, $\alpha\in(0,1/2)$, and $\varepsilon>0$
such that the following holds. Let $r_d$ be an odd integer with
$r_d/d\to\alpha$, and set
\[
 n_d=\left\lceil A D_{r_d}\log(D_{r_d}+1)\right\rceil.
\]
For all sufficiently large $d$, the Gaussian sign posterior trained on
$n_d$ independent uniform parity examples satisfies
\[
 \mathbb P_S\{F_{13/20}(S)\ge\varepsilon\}\ge\frac34.
\]
In particular $n_d=\exp((h(\alpha)+o(1))d)$, where
$h(\alpha)=-\alpha\log\alpha-(1-\alpha)\log(1-\alpha)>0$.
The same conclusion holds with $F_{13/20}(S)$ replaced by
$\liminf_{N\to\infty}F_N(S)$ in any finite-width model whose conditional
prediction probabilities converge to this Gaussian sign posterior.
This includes the negative-$\gamma$ scaling width limit.
\end{theorem}
\begin{proof}
Fix $B=2$ and $\varepsilon=c^2/(1024B)$. Set
$u=c^2/(2048B)$ and $\delta=1/20$. First choose $A$ large enough
for Lemma~\ref{lem:schedule-relu} with $\eta=1/4$ and this $\delta$.
The class $\{x: f(x)p(x)<0,\ p\in V_r\}$ has VC dimension at most
$D_r+1$. Thus the uniform VC theorem \cite{vapnik1971} gives
\eqref{eq:schedule-vc} with probability at least $1-\delta$ for all
large $d$, since $n_d/D_{r_d}\to\infty$.

The spectral estimates give, for odd $r$ and $r/d$ sufficiently small,
\[
 \frac{n_d\lambda_{r+1}}{\tau_r}
 \le C_L A\alpha^2\log(1/\alpha)+o(1).
\]
Corollary~\ref{cor:gram-schedule} gives
$\mathbb P_S\{\|Q_{SS}\|_{\mathrm{op}}>2\tau_r\}\le\delta$
for all large $d$ after $\alpha$ is made sufficiently small, depending
on $A,L,\delta$. Decrease $\alpha$ further so that the preceding spectral
ratio also meets Proposition~\ref{prop:schedule-deterministic}.
A union bound over the three sampling events proves a probability of
at least $1-3\delta>3/4$.

For the final assertion, fix $d$ and the sampled training set. At every
input with limiting correct-label probability at most $13/20$, the
finite-width correct-label probability is eventually below $2/3$.
The cube is finite, so this holds simultaneously at every such input,
which proves $\liminf_N F_N(S)\ge F_{13/20}(S)$.
\end{proof}

The theorem concerns the displayed sample schedule. It does not assert
failure at every smaller sample size, and no monotonicity of posterior
confidence in sample size is assumed. The later paper~\cite{gaussian-all-n} covers all intervening sample sizes
by an additional middle-band argument. The proof
also uses convexity of Gaussian sign conditioning; it does not establish
the analogous exponential conclusion at $\gamma=0$ or $0<\gamma<1$.


\begin{thebibliography}{9}
\bibitem{gaussian-all-n}
ChatGPT-6 Astra, K. H\"anni, and D. Vaintrob.
\emph{Exponential sample lower bounds for Bayesian Gaussian process
classification}. Companion manuscript, October 2026.
\bibitem{subcritical-exponential}
ChatGPT-6 Astra, K. H\"anni, and D. Vaintrob.
\emph{Exponential parity sample complexity below the mean-field scaling}.
Companion manuscript, October 2026.

\bibitem{flajolet2009}
P. Flajolet and R. Sedgewick.
\emph{Analytic Combinatorics}. Cambridge University Press, 2009.
Chapter VI, algebraic singularity analysis and multiple dominant
singularities.
\bibitem{mai2021coresets}
T. Mai, C. Musco, and A. Rao.
Coresets for Classification---Simplified and Strengthened.
\emph{Advances in Neural Information Processing Systems} 34, 2021.
\href{https://papers.neurips.cc/paper_files/paper/2021/file/6098ed616e715171f0dabad60a8e5197-Paper.pdf}{Primary paper}; Theorem 5.
\bibitem{tropp2008}
J. A. Tropp.
Norms of random submatrices and sparse approximation.
\emph{Comptes Rendus Math\'ematique} 346:1271--1274, 2008.
\href{https://tropp.caltech.edu/papers/Tro08-Norms-Random-v3.pdf}{Author manuscript}; Theorem 1.1.
\bibitem{vapnik1971}
V. N. Vapnik and A. Ya. Chervonenkis.
On the uniform convergence of relative frequencies of events to their probabilities.
\emph{Theory of Probability and its Applications} 16:264--280, 1971.
\bibitem{prekopa1973}
A. Pr\'ekopa.
On logarithmic concave measures and functions.
\emph{Acta Scientiarum Mathematicarum} 34:335--343, 1973.
\end{thebibliography}
\end{document}
